\documentclass[10pt]{article}

% ───────────────────────── Paquetes ─────────────────────────
\usepackage[utf8]{inputenc}
\usepackage[T1]{fontenc}
\usepackage[a4paper,top=1.6cm,bottom=1.8cm,left=1.6cm,right=1.6cm]{geometry}
\usepackage{helvet}
\renewcommand{\familydefault}{\sfdefault}
\usepackage{xcolor}
\usepackage{colortbl}
\usepackage{array}
\usepackage{tabularx}
\usepackage{longtable}
\usepackage{booktabs}
\usepackage{pdflscape}
\usepackage{ragged2e}
\usepackage[most]{tcolorbox}
\usepackage{enumitem}
\usepackage{fancyhdr}
\usepackage{lastpage}
\usepackage{multirow}

% ───────────────────────── Colores UTEQ ─────────────────────────
\definecolor{uteqverde}{HTML}{1D6B36}
\definecolor{uteqazul}{HTML}{1A5276}
\definecolor{uteqoro}{HTML}{B8860B}
\definecolor{filaazul}{HTML}{EAF2F8}
\definecolor{filaoro}{HTML}{FCF3DC}
\definecolor{filaverde}{HTML}{E8F4EC}
\definecolor{grisclaro}{HTML}{F4F6F7}
\definecolor{grisborde}{HTML}{BDC3C7}
\definecolor{rojoadm}{HTML}{922B21}

% ───────────────────────── Tipos de columna ─────────────────────────
\newcolumntype{L}[1]{>{\raggedright\arraybackslash}p{#1}}
\newcolumntype{C}[1]{>{\centering\arraybackslash}p{#1}}

% ───────────────────────── Encabezado / pie ─────────────────────────
\pagestyle{fancy}
\fancyhf{}
\renewcommand{\headrulewidth}{0pt}
\fancyfoot[L]{\footnotesize\color{uteqazul}ISR-401 Ingenier\'ia de Requisitos · 2026--2027 PPA}
\fancyfoot[C]{}
\fancyfoot[R]{\footnotesize\color{uteqazul}P\'ag.\ \thepage/\pageref{LastPage}}

\setlength{\parindent}{0pt}
\renewcommand{\arraystretch}{1.25}

% ───────────────────────── Comandos auxiliares ─────────────────────────
\newcommand{\nivelhead}[2]{\cellcolor{#1}\textcolor{white}{\textbf{#2}}}

\begin{document}

% ════════════════════ TÍTULO ════════════════════
\begin{tcolorbox}[colback=uteqverde,colframe=uteqverde,sharp corners,boxrule=0pt,
  left=10pt,right=10pt,top=8pt,bottom=8pt]
  \begin{center}
    {\large\textbf{\textcolor{white}{UNIVERSIDAD T\'ECNICA ESTATAL DE QUEVEDO}}}\\[2pt]
    {\small\textcolor{white}{Facultad de Ciencias de la Computaci\'on y Dise\~no Digital — Carrera de Ingenier\'ia en Software}}\\[6pt]
    {\Large\textbf{\textcolor{white}{R\'UBRICA DE EVALUACI\'ON — EXPOSICI\'ON}}}\\[3pt]
    {\normalsize\textcolor{white}{«Herramientas Automatizadas para la Ingenier\'ia de Requisitos» (Herramientas CASE — Segundo Corte)}}
  \end{center}
\end{tcolorbox}

\vspace{4pt}

% ════════════════════ DATOS GENERALES ════════════════════
\noindent\begin{tabularx}{\linewidth}{|L{0.16\linewidth}|X|L{0.16\linewidth}|X|}
\hline
\cellcolor{filaazul}\textbf{Asignatura} & ISR-401 Ingenier\'ia de Requisitos &
\cellcolor{filaazul}\textbf{Per\'iodo} & 2026--2027 PPA \\
\hline
\cellcolor{filaazul}\textbf{Docente} & Guerrero Ulloa Gleiston Cicer\'on &
\cellcolor{filaazul}\textbf{Paralelo} & \rule{2.5cm}{0.4pt} \\
\hline
\cellcolor{filaazul}\textbf{Equipo / N.\textsuperscript{o}} & \rule{4cm}{0.4pt} &
\cellcolor{filaazul}\textbf{Fecha} & \rule{2.5cm}{0.4pt} \\
\hline
\cellcolor{filaazul}\textbf{Integrantes} & \multicolumn{3}{X|}{\rule{0pt}{12pt}} \\
\hline
\end{tabularx}

\vspace{8pt}

% ════════════════════ CRITERIOS DE ADMISIÓN ════════════════════
\begin{tcolorbox}[colback=white,colframe=rojoadm,sharp corners,boxrule=1pt,
  left=8pt,right=8pt,top=6pt,bottom=6pt,
  title={\textbf{CRITERIOS DE ADMISI\'ON (\textit{gatekeeper})}},
  coltitle=white,colbacktitle=rojoadm,fonttitle=\bfseries]
  El incumplimiento de \textbf{cualquiera} de los siguientes criterios fija la nota en el
  \textbf{m\'inimo} (10\,\% del m\'aximo) y la r\'ubrica no se eval\'ua.
  \begin{itemize}[leftmargin=1.4em,itemsep=1pt,topsep=3pt]
    \item[\textbf{A1.}] Entrega del informe escrito \textbf{y} de las diapositivas.
    \item[\textbf{A2.}] Identificaci\'on completa del equipo (integrantes, curso/paralelo y docente).
    \item[\textbf{A3.}] An\'alisis de al menos \textbf{dos (2)} herramientas automatizadas de IR.
    \item[\textbf{A4.}] Exposici\'on realizada en la fecha y turno asignados.
  \end{itemize}
\end{tcolorbox}

\vspace{6pt}

% ════════════════════ ESCALA DE VALORACIÓN ════════════════════
\noindent\textbf{\textcolor{uteqazul}{Escala de valoraci\'on por criterio}}\\[3pt]
\noindent\begin{tabularx}{\linewidth}{|C{0.10\linewidth}|L{0.24\linewidth}|X|}
\hline
\rowcolor{uteqazul}\textcolor{white}{\textbf{Nivel}} & \textcolor{white}{\textbf{Denominaci\'on}} & \textcolor{white}{\textbf{Valor}} \\
\hline
\cellcolor{filaverde}\textbf{N4} & Excelente & 4 — equivale al 100\,\% del peso del criterio \\
\hline
\cellcolor{filaazul}\textbf{N3} & Satisfactorio & 3 — equivale al 75\,\% del peso del criterio \\
\hline
\cellcolor{filaoro}\textbf{N2} & En desarrollo & 2 — equivale al 50\,\% del peso del criterio \\
\hline
\cellcolor{grisclaro}\textbf{N1} & Insuficiente & 1 — equivale al 25\,\% del peso del criterio \\
\hline
\end{tabularx}

\vspace{10pt}
\begin{center}
\textit{\small La tabla anal\'itica de criterios contin\'ua en la siguiente p\'agina (orientaci\'on horizontal).}
\end{center}

% ════════════════════ TABLA ANALÍTICA (LANDSCAPE) ════════════════════
\begin{landscape}
\thispagestyle{fancy}
\begin{center}
{\large\textbf{\textcolor{uteqazul}{Tabla anal\'itica de criterios de evaluaci\'on}}}
\end{center}
\vspace{2pt}
\footnotesize
\renewcommand{\arraystretch}{1.18}
\setlength{\tabcolsep}{4pt}

\begin{longtable}{|L{3.0cm}|L{5.7cm}|L{5.4cm}|L{4.9cm}|L{4.3cm}|}
\hline
\rowcolor{uteqazul}
\textcolor{white}{\textbf{Criterio (peso)}} &
\nivelhead{uteqverde}{N4 — Excelente (100\%)} &
\nivelhead{uteqazul}{N3 — Satisfactorio (75\%)} &
\nivelhead{uteqoro}{N2 — En desarrollo (50\%)} &
\nivelhead{rojoadm}{N1 — Insuficiente (25\%)} \\
\hline
\endfirsthead
\hline
\rowcolor{uteqazul}
\textcolor{white}{\textbf{Criterio (peso)}} &
\nivelhead{uteqverde}{N4 — Excelente} &
\nivelhead{uteqazul}{N3 — Satisfactorio} &
\nivelhead{uteqoro}{N2 — En desarrollo} &
\nivelhead{rojoadm}{N1 — Insuficiente} \\
\hline
\endhead

% ---------- C1 ----------
\cellcolor{filaazul}\textbf{C1. Dominio del contenido t\'ecnico}\newline\textit{(20\%)} &
Describe cada herramienta con precisi\'on conceptual; funcionalidades fundamentadas en literatura (Pohl 2010, IEEE 29148:2018, SWEBOK v4, art\'iculos indexados); terminolog\'ia de IR correcta y uniforme. &
Contenido correcto y fundamentado, con profundidad desigual entre herramientas o errores menores de terminolog\'ia/redacci\'on. &
Descripci\'on superficial tipo folleto comercial; escasa fundamentaci\'on acad\'emica. &
Contenido err\'oneo, copiado o sin fundamentaci\'on. \\
\hline

% ---------- C2 ----------
\cellcolor{filaazul}\textbf{C2. Comparativa de funcionalidades}\newline\textit{(15\%)} &
Tabla comparativa con $\geq 5$ criterios pertinentes (trazabilidad, gesti\'on de cambios, colaboraci\'on, IA, costo, facilidad de uso), con an\'alisis interpretativo y conclusi\'on fundamentada. &
Tabla completa pero an\'alisis interpretativo breve o parcial. &
Tabla con criterios insuficientes o sin an\'alisis. &
Sin comparativa. \\
\hline

% ---------- C3 ----------
\cellcolor{filaazul}\textbf{C3. Ventajas y limitaciones basadas en evidencia}\newline\textit{(15\%)} &
Ventajas y limitaciones de \textbf{todas} las herramientas, cada afirmaci\'on respaldada con cita verificable; sin juicios gratuitos. &
Cobertura completa pero con afirmaciones parcialmente sin respaldo. &
Listas gen\'ericas sin citas o incompletas para alguna herramienta. &
Ausentes o triviales. \\
\hline

% ---------- C4 ----------
\cellcolor{filaazul}\textbf{C4. Aplicabilidad al Proyecto Fin de Curso (PFC)}\newline\textit{(15\%)} &
Vincula funcionalidades concretas de cada herramienta con requisitos y artefactos \textbf{reales} del PFC (m\'odulos, RNF, normativa del dominio), indicando en qu\'e actividad del proceso de IR aporta. &
Vinculaci\'on concreta pero solo para parte de las herramientas o sin mapear a actividades del proceso de IR. &
Aplicabilidad gen\'erica, v\'alida para cualquier proyecto. &
Ausente. \\
\hline

% ---------- C5 ----------
\cellcolor{filaazul}\textbf{C5. Manual de usuario con capturas de pantalla}\newline\textit{(20\%)} &
El documento incluye una secci\'on tipo manual de usuario: pasos de instalaci\'on/acceso y uso de al menos una herramienta, con capturas \textbf{propias} del equipo (datos/requisitos del PFC visibles), numeradas como figuras y referenciadas en el texto. &
Manual con capturas propias pero pasos incompletos o capturas sin numerar/referenciar. &
Solo capturas aisladas o im\'agenes del sitio oficial/stock, sin secuencia de pasos ni evidencia de uso propio. &
Sin manual ni capturas de uso. \\
\hline

% ---------- C6 ----------
\cellcolor{filaazul}\textbf{C6. Material de exposici\'on (diapositivas)}\newline\textit{(10\%)} &
Dise\~no consistente y profesional; jerarqu\'ia visual clara; texto sint\'etico; tablas legibles; im\'agenes licenciadas/propias sin marcas de agua; secuencia l\'ogica completa. &
Dise\~no correcto con sobrecarga de texto o desbalance puntual. &
Inconsistencias notorias, im\'agenes con marca de agua, texto duplicado o ilegible. &
Diapositivas improvisadas o ausentes. \\
\hline

% ---------- C7 ----------
\cellcolor{filaazul}\textbf{C7. Referencias y rigor acad\'emico}\newline\textit{(5\%)} &
Formato IEEE uniforme; todas las referencias verificables y con datos internamente consistentes (autores–t\'itulo–\textit{venue}–DOI); sin duplicados; todas citadas en el texto. &
IEEE correcto con fallas puntuales (duplicados, 1--2 entradas incompletas o no citadas). &
M\'ultiples referencias incompletas, no verificables o no citadas. &
Sin referencias o no verificables en su mayor\'ia. \\
\hline

\end{longtable}
\end{landscape}

% ════════════════════ HOJA DE CÁLCULO ════════════════════
\noindent\textbf{\textcolor{uteqazul}{Hoja de c\'alculo de la nota}}\\[4pt]
\noindent\begin{tabularx}{\linewidth}{|C{0.07\linewidth}|X|C{0.12\linewidth}|C{0.14\linewidth}|C{0.18\linewidth}|}
\hline
\rowcolor{uteqazul}
\textcolor{white}{\textbf{Crit.}} & \textcolor{white}{\textbf{Descripci\'on}} &
\textcolor{white}{\textbf{Peso}} & \textcolor{white}{\textbf{Nivel (1--4)}} &
\textcolor{white}{\textbf{Aporte (Nivel\,$\times$\,Peso)}} \\
\hline
C1 & Dominio del contenido t\'ecnico & 0.20 & & \\ \hline
\rowcolor{grisclaro} C2 & Comparativa de funcionalidades & 0.15 & & \\ \hline
C3 & Ventajas y limitaciones basadas en evidencia & 0.15 & & \\ \hline
\rowcolor{grisclaro} C4 & Aplicabilidad al PFC & 0.15 & & \\ \hline
C5 & Manual de usuario con capturas de pantalla & 0.20 & & \\ \hline
\rowcolor{grisclaro} C6 & Material de exposici\'on (diapositivas) & 0.10 & & \\ \hline
C7 & Referencias y rigor acad\'emico & 0.05 & & \\ \hline
\multicolumn{2}{|r|}{\cellcolor{filaazul}\textbf{Suma de aportes}} & \cellcolor{filaazul}\textbf{1.00} & &
\cellcolor{filaazul}\rule{0pt}{12pt} \\ \hline
\end{tabularx}

\vspace{8pt}

% ════════════════════ FÓRMULA ════════════════════
\begin{tcolorbox}[colback=filaverde,colframe=uteqverde,sharp corners,boxrule=0.8pt,
  left=8pt,right=8pt,top=6pt,bottom=6pt]
  \textbf{F\'ormula de c\'alculo}
  \[
  \text{Nota} \;=\; \left( \frac{\displaystyle\sum_{i=1}^{7} \text{Nivel}_i \times \text{Peso}_i}{4} \right) \times E
  \]
  donde $\text{Nivel}_i \in \{1,2,3,4\}$, $\text{Peso}_i$ es la fracci\'on del criterio ($\sum \text{Peso}_i = 1$)
  y $E$ es el factor de escala:
  \begin{itemize}[leftmargin=1.4em,itemsep=1pt,topsep=2pt]
    \item $E = 10$ \;\; escala decimal (0--10).
    \item $E = 3$ \;\; exposici\'on del Segundo Corte (m\'aximo \textbf{3.00} puntos).
  \end{itemize}
  \textit{Piso por gatekeeper:} si se incumple un criterio de admisi\'on, $\text{Nota} = 0.10 \times E$
  (1.0 en escala 10; 0.30 en escala 3).
\end{tcolorbox}

\vspace{6pt}

% ════════════════════ INTERPRETACIÓN ════════════════════
\noindent\textbf{\textcolor{uteqazul}{Interpretaci\'on de la nota}}\\[4pt]
\noindent\begin{tabularx}{\linewidth}{|L{0.32\linewidth}|C{0.28\linewidth}|>{\centering\arraybackslash}X|}
\hline
\rowcolor{uteqazul}
\textcolor{white}{\textbf{Desempe\~no}} & \textcolor{white}{\textbf{\% del m\'aximo}} & \textcolor{white}{\textbf{Equivalente en escala 3.0}} \\
\hline
\cellcolor{filaverde}Excelente & $\geq 90\,\%$ & $\geq 2.70$ \\ \hline
\cellcolor{filaazul}Satisfactorio & $70\%$ -- $89\%$ & $2.10$ -- $2.69$ \\ \hline
\cellcolor{filaoro}En desarrollo & $50\%$ -- $69\%$ & $1.50$ -- $2.09$ \\ \hline
\cellcolor{grisclaro}Insuficiente & $< 50\,\%$ & $< 1.50$ \\ \hline
\end{tabularx}

\vspace{10pt}

% ════════════════════ OBSERVACIONES / FIRMA ════════════════════
\noindent\begin{tabularx}{\linewidth}{|X|}
\hline
\cellcolor{filaazul}\textbf{Observaciones / retroalimentaci\'on} \\
\hline
\rule{0pt}{40pt} \\
\hline
\end{tabularx}

\vspace{18pt}
\noindent\begin{tabularx}{\linewidth}{C{0.5\linewidth}C{0.5\linewidth}}
\rule{6cm}{0.4pt} & \rule{6cm}{0.4pt} \\
Nota final (sobre 3.00) & Firma del docente \\
\end{tabularx}

\end{document}
